%%!TEX root = the-subsets.tex
\chapter{Subsets}\label{chapter:shefel}

Here we give a partial answer to the following question:
\textit{Which subsets of Euclidean space, equipped with their induced length-metrics, are  $\CAT(0)$?}

\section{Motivating examples}\label{sec:solid-parabolonds}

Consider three subgraphs of different quadric surfaces:
\begin{align*}
A&=\set{(x,y,z)\in\EE^3}{z\le x^2+y^2},
\\
B&=\set{(x,y,z)\in\EE^3}{z\le -x^2-y^2},
\\
C&=\set{(x,y,z)\in\EE^3}{z\le x^2-y^2}.
\end{align*}

\begin{thm}{Question}\label{CAT(0)?}
Which of the sets $A$, $B$ and $C$, if equipped with the induced length metric, are $\CAT(0)$ and why?
\end{thm}

The answers are given below, but it is instructive to try to answer yourself before reading further.

\parit{$\bm{A}$.}
No, $A$ is not $\CAT(0)$.

Choose a triangle $[pqr]$ in the $(x,y)$-plane that surrounds the origin.
Lift it a bit up, so $[pqr]$ remains in $A$, and observe that this triangle is not thin in the induced length metric of $A$.
\qeds

Let us give another visual argument more in the style of  differential geometry.
 
\parit{Another way.} The boundary $\partial A$ is the paraboloid described by  $z\z=x^2+y^2$;  in particular it bounds an open convex set in $\EE^3$ whose complement is~$A$.
The nearest-point projection of $A\to\partial A$ is short (Exercise~\ref{ex:closest-point}).
It follows that $\partial A$ is a convex set in $A$ equipped with its induced length metric.

Therefore if $A$ is $\CAT(0)$, then so is~$\partial A$.
The latter is not true: $\partial A$ is a smooth convex surface, and has strictly positive curvature by the Gauss formula.
\qeds

\parit{$\bm{B}$.} Yes, $B$ is $\CAT(0)$.

Evidently $B$ is a closed convex set in~$\EE^3$.
Therefore the length metric on $B$ coincides with the Euclidean metric
and $\CAT(0)$ comparison holds.\qeds

\parit{$\bm{C}$.} Yes, $C$ is $\CAT(0)$,
but the proof is not as easy as before.

Set $f_t(x,y)=x^2-y^2 -2\cdot (x-t)^2$.
Consider the one-parameter family of sets 
\[V_t=\set{(x,y,z)\in\EE^3}{z\le f_t(x,y)}.\]

\begin{figure}[t!]
\vskip-10mm
\centering
\includegraphics{mppics/pic-1201}
\caption{$C=\set{(x,y,z)\in\EE^3}{z\le x^2-y^2}$.}
\end{figure}

Each set $V_t$ is a solid paraboloid tangent to $\partial C$ along the parabola $y\z\mapsto(t,y,t^2-y^2)$.
The set $V_t$ is closed and convex for any $t$, and
\[C=\bigcup_t V_t.\]
Further note that the function $t\mapsto f_t(x,y)$ is concave for any fixed $x,y$.
Therefore
\begin{clm}{}\label{eq:VcVnV-0}
if $a<b<c$, then $V_b\supset V_a\cap V_c$.
\end{clm}

Given $t_1,\ldots, t_n\in \RR$ consider the  union
\[C'=V_{t_1}\cup\ldots\cup V_{t_n}.\]
The inclusion \ref{eq:VcVnV-0} makes it possible to apply Reshetnyak gluing theorem \ref{thm:gluing} recursively and show that $C'$ is $\CAT(0)$.

By approximation, the $\CAT(0)$ comparison holds for any 4 points in $C$;
hence $C$ is $\CAT(0)$.
More precisely, choose $x_1,x_2,x_3,x_4\in C$ and 6 geodesics $[x_ix_j]_C$ between them.
Choose $\eps>0$, shift each $[x_ix_j]_C$ down by $\eps$, and reconnect it to $x_i$ and $x_j$ by vertical $\eps$-segments.
Denote the obtained curve by $\gamma_{i,j}$;
note that 
\[\length\gamma_{i,j}=\dist{x_i}{x_j}{C}+2\cdot\eps.\]
We may assume that $C'$ contains all $\gamma_{i,j}$.
It follows that 
\[\dist{x_i}{x_j}{C}\le \dist{x_i}{x_j}{C'}\le \dist{x_i}{x_j}{C}+2\cdot\eps\]
Since $\eps>0$ is arbitrary and $C'$ is $\CAT(0)$, so is $C$.\qeds

\parbf{Remarks.}
The last argument will be reused in the proof of \ref{thm:set-with-smooth-bry:CBA}.
While the set $C$ is not convex, we shall see that it is \textit{two-convex} as defined in the next section.
In fact, two-convexity is closely related to the inheritance of an upper curvature bound by a subset.

\section{Two-convexity}

\begin{thm}{Definition}\label{def:two-convex}
We say that a subset $K\subset \EE^m$ is \index{two-convex set}\emph{two-convex}
if the following condition holds for any plane $W\subset\EE^m$:
If $\gamma$ is a simple closed curve in $W\cap K$ 
that is null-homotopic in $K$,  
then it is null-homotopic in $W\cap K$, and in particular the disc in $W$ bounded by $\gamma$ lies in~$K$.
\end{thm}

Note that two-convex sets do not have to be connected or simply connected. 
The following two propositions follow immediately from the definition.

\begin{thm}{Proposition}
Any subset in $\EE^2$ is two-convex.
\end{thm}

\begin{thm}{Proposition}\label{prop:two-hull}
The intersection of an arbitrary collection of two-convex sets in $\EE^m$ is two-convex.
\end{thm}

\begin{thm}{Proposition}\label{prop:two-hull-open}
The interior of any two-convex set in $\EE^m$ is a two-convex set.
\end{thm}

\parit{Proof.}
Fix a  two-convex set $K\subset \EE^m$ and a 2-plane $W$;
denote by \index{3@$\Int$ (interior)}$\Int K$ the interior of $K$.
Let $\gamma$ be a closed simple curve in $W\cap \Int K$ 
that  is contractible in the interior of~$K$.

Since $K$ is two-convex,
the plane disc $D$ bounded by $\gamma$ lies in~$K$.
The same holds for the translations of $D$ by small vectors.
Therefore $D$ lies in $\Int K$;
that is, $\Int K$ is two-convex.
\qeds

\begin{thm}{Definition}
Given a subset $K\subset \EE^m$, define its two-convex hull (briefly, $\Conv_2K$) as the intersection of all two-convex subsets containing~$K$.
\end{thm}

Note that by Proposition~\ref{prop:two-hull},
the two-convex hull of any set is two-convex.
Further, 
by \ref{prop:two-hull-open}, the
two-convex hull of an open set is open.

The next proposition describes  closed two-convex sets with smooth boundary.

\begin{thm}{Proposition}\label{prop:two-cove+smooth}
Let $K\subset \EE^m$ be a closed subset.

Assume that the boundary of $K$ is a smooth hypersurface~$S$.
Consider the unit normal vector field $\nu$ on $S$ that  points outside of~$K$.
Denote by $k_1\le \ldots\le k_{m-1}$ the principal curvature functions of $S$ with respect to $\nu$ (note that if $K$ is convex, then $k_1\ge 0$, and for connected $K$ the converse is also true).

Then $K$ is two-convex if and only if $k_2(p)\ge 0$ for any point $p\in S$.
Moreover, if $k_2(p)<0$ at some point $p$, then Definition~\ref{def:two-convex} fails for some curve $\gamma$ forming a triangle in an arbitrarily small neighborhood of~$p$.
\end{thm}

The proof uses a straightforward modification of the Morse theory for manifolds with boundary; the paper of Sergei Vakhrameev \cite{vakhrameev} contains all the necessary lemmas.

\begin{figure}[t!]
\vskip-0mm
\centering
\includegraphics{mppics/pic-1202}
\caption{A triangle moved in the direction of $\nu(p)$.}
\end{figure}

\parit{Proof; only-if part.}
Suppose $k_2(p)<0$ at some point $p\in S$,
consider the plane $W$ containing $p$ and spanned by the first two principal directions at~$p$.
Choose a small triangle in $W$ that surrounds $p$ and move it slightly in the direction of $\nu(p)$.
We get a triangle $\trig xyz$ that is null-homotopic in $K$,
but the solid triangle $\Delta=\Conv\{x,y,z\}$ bounded by $\trig xyz$ does not lie in $K$ completely.
Therefore $K$ is not two-convex.

\parit{If part.}
Recall that a smooth function $f\:\EE^m\to\RR$ is called \index{strongly convex function}\emph{strongly convex} if its Hessian is positive definite at each point.

Suppose $f\:\EE^m\to\RR$ is a smooth
 strongly
 convex function such that the restriction $f|_S$ is a Morse function.
Note that a generic smooth 
strongly 
convex function $f\:\EE^m\to\RR$ has this property.

For a critical point $p$ of $f|_S$, the outer normal vector $\nu(p)$ is parallel to the gradient $\nabla_pf$;
we say that $p$ is a 
\index{positive critical point}\emph{positive critical point}
if $\nu(p)$ and $\nabla_p f$ point in the same direction, 
and 
\index{negative critical point}\emph{negative} otherwise.
If $f$ is generic, then we can assume that the sign is defined for all critical points;
that is, $\nabla_pf\ne0$ for any critical point $p$ of~$f|_S$.

Since $k_2\ge 0$ and the function $f$ is  strongly
 convex, 
the negative critical points of $f|_S$
have index at most~1.

Given a real value $s$, set 
\[K_s=\set{x\in K}{f(x)<s}.\]
Assume  $\phi_0\:\DD\to K$ is a continuous map of the disc $\DD$
such that $\phi_0(\partial \DD)\subset K_s$.

Note that by the Morse lemma, 
there is a homotopy $\phi_t\:\DD\to K$ rel $\partial \DD$ such that 
$\phi_1(\DD)\subset K_s$.

Indeed, we can construct a homotopy $\phi_t\:\DD\to K$ that decreases the maximum of $f\circ\phi$ on $\DD$ until the maximum occurs at a critical point $p$ of~$f|_S$.
This point cannot be negative; otherwise, its index would be at least 2.
If this critical point is positive, then it is easy to decrease the maximum a little by pushing the disc from $S$ into $K$ in the direction of $-\nabla f_p$.

Consider a closed curve $\gamma\:\mathbb{S}^1\to K$ that is null-homotopic in~$K$.
Note that the distance function \[f_0(x)\z=\dist{\Conv\gamma}{x}{\EE^m}\] 
is convex on $\EE^m$ by Exercise \ref{ex:displacement}.
Therefore $f_0$ can be approximated by smooth strongly convex functions $f$ in general position.
From above, there is a disc in $K$ with boundary $\gamma$
that lies arbitrarily close to $\Conv\gamma$.
Since $K$ is closed, the statement follows.
\qeds

Note that the ``if'' part proves a somewhat stronger statement.
Namely,  any plane curve $\gamma$ (not necessary simple) which is  contractible in $K$ is also contractible in the intersection of $K$ with the plane of~$\gamma$.
The latter condition does not hold for the complement 
of two planes in $\EE^4$, which is two-convex by Proposition~\ref{prop:two-hull};
see also Exercise~\ref{ex:two-planes} below.
By the following proposition, there are no such examples in~$\EE^3$.

\begin{thm}{Proposition}\label{prop:3d-strong-2-convexity}
Let $\Omega\subset\EE^3$ be an open two-convex subset.
Then for any plane $W\subset\EE^3$, 
any closed curve in $W\cap \Omega$ 
that is null-homotopic in $\Omega$ is also null-homotopic in $W\cap \Omega$.
\end{thm}

This statement is intuitively obvious, but the proof is not trivial;
it uses the following classical result.

\begin{thm}{Loop theorem}
Let $M$ be a 3-dimensional manifold with nonempty boundary $\partial M$.
Assume 
$f\: (\DD,\partial \DD)\to (M,\partial M)$
is a continuous map of pairs
such that the boundary curve $f|_{\partial \DD}$ is not null-homotopic in $\partial M$; here $\DD$ denotes the disc.
Then there is an embedding of pairs $h\: (\DD,\partial \DD)\to (M,\partial M)$ with the same property.
\end{thm}

The theorem is due to Christos Papakyriakopoulos \cite[a proof can be found in][]{hatcher}. 

\parit{Proof of \ref{prop:3d-strong-2-convexity}.}
Fix a closed plane curve $\gamma$ in $W\cap \Omega$ that  is null-homotopic in $\Omega$. Suppose $\gamma$ is not contractible in  $W\cap \Omega$.

Let $\phi\: \DD\to\Omega$ be a map of the disc with the boundary curve~$\gamma$.

Since $\Omega$ is open we can first change  $\phi$  slightly so that $\phi(x)\notin W$ for $1-\eps\z<|x|\z<1$ for some small $\eps>0$.
By further changing $\phi$ slightly we can assume that it is transversal to $W$  on  $\Int \DD$ and agrees with the previous map near~$\partial \DD$.

This means that $\phi^{-1} (W)\cap \Int \DD$ consists of finitely many simple closed curves which cut $\DD$ into several components. 
Consider one of the ``innermost'' components $c'$;
that is, $c'$ is a boundary curve of a disc $\DD'\subset \DD$,
$\phi(c')$ is a closed curve in $W$ and $\phi(\DD')$  completely lies in one of the two half-spaces  with boundary~$W$. 
Denote this half-space by~$H$.

If $\phi(c')$ is not contractible in $W\cap \Omega$, then applying the loop theorem to $M^3=H\cap \Omega$ we conclude that  there exists a {}\emph{simple} closed curve $\gamma'\subset \Omega\cap W$ which is not contractible in $\Omega\cap W$ but is contractible in $\Omega\cap H$. 
This contradicts two-convexity of~$\Omega$. 

Hence $\phi(c')$ is contractible in $W\cap \Omega$. Therefore $\phi$ can be changed in a small neighborhood $U$ of $\DD'$ so that the new map $\hat\phi$ maps $U$ to one side of~$W$. 
In particular, the set $\hat\phi^{-1}(W)$ consists of the same curves as $\phi^{-1} (W)$ with the exception of~$c'$.

Repeating this process several times we reduce the problem to the case where $\phi^{-1} (W)\cap \Int \DD=\emptyset$.
This means that $\phi(\DD)$ lies entirely in one of the half-spaces bounded by $W$.

Again applying the loop theorem, we obtain a simple closed curve in $W\cap \Omega$ which is not contractible in $W\cap \Omega$ but is contractible in~$\Omega$. 
This again contradicts two-convexity of~$\Omega$. 
Hence $\gamma$ is contractible in  $W\cap \Omega$ as claimed.
\qeds

\section{Sets with smooth boundary}\label{sec:smooth-bry}

In this section, we characterize the subsets with smooth boundary in $\EE^m$ that form $\CAT(0)$ spaces. 

\begin{thm}{Smooth two-convexity theorem}\label{thm:set-with-smooth-bry:CBA}
Let $K$ be a closed, simply connected subset in $\EE^m$ equipped with the induced length metric.
Assume $K$ is bounded by a smooth hypersurface.
Then 
$K$ is $\CAT(0)$ if and only if $K$ is two-convex.
\end{thm}

The proof below is based on the argument in \ref{CAT(0)?}$C$.

\parit{Proof.}
Denote by $S$ and by $\Omega$ the boundary and the interior of $K$ respectively. 
Since $K$ is connected and $S$ is smooth, $\Omega$ is also connected.

Denote by $k_1(p)\le\ldots\le k_{m-1}(p)$ the principal curvatures of $S$ at $p\in S$ with respect to the normal vector $\nu(p)$ pointing out of~$K$.
By Proposition~\ref{prop:two-cove+smooth}, $K$ is two-convex if and only if $k_2(p)\ge 0$ for any $p\in S$.

\parit{Only-if part.}
Assume $K$ is not two-convex.
Then by Proposition~\ref{prop:two-cove+smooth}, $k_2(p)<0$ at some point $p\in S$. Hence there is a triangle $\trig xyz$ in $K$ that is null-homotopic in $K$, but the solid triangle $\Delta=\Conv\{x,y,z\}$ does not lie in $K$ completely.
Evidently the triangle $\trig xyz$ is not thin in~$K$. 
Hence $K$ is not $\CAT(0)$.

\parit{If part.}
Since $K$ is simply connected,
by the globalization theorem (\ref{thm:hadamard-cartan})
it suffices to show that any point $p\in K$ admits a $\CAT(0)$ neighborhood.

Evidently, such a neighborhood exists if $p\in\Int K$

Assume $p\in S$ and $k_2(p)>0$.
Fix a sufficiently small $\eps>0$ and set $K'\z=K\cap \cBall[p,\eps]$.
Let us show that 
\begin{clm}{}\label{K'-is-CAT}
$K'$ is $\CAT(0)$.
\end{clm}

Consider the coordinate system with the origin at $p$
and the principal directions and $\nu(p)$ as the coordinate directions.
For small $\eps>0$, the set $K'$ 
can be described as a subgraph
\[K'
=
\set
{(x_1,\ldots,x_m)\in \cBall[p,\eps]}
{x_m\le f(x_1,\ldots,x_{m-1})}.\]

Fix $s\in[-\eps,\eps]$.
Since $\eps$ is small and $k_2(p)>0$, the restriction 
$f|_{x_1=s}$
is concave in the $(m-2)$-dimensional cube defined by the inequalities $|x_i|<2\cdot\eps$ for $2\le i\le m-1$.

Fix a negative real value $\lambda<k_1(p)$.
Given $s\in (-\eps,\eps)$,
consider the set 
\[V_s
=
\set
{(x_1,\ldots,x_m)\in K'}
{x_m\le f(x_1,\ldots,x_{m-1})+\lambda\cdot (x_1-s)^2}.\]
Note that the function 
\[(x_1,\ldots, x_{m-1})\mapsto f(x_1,\ldots,x_{m-1})+\lambda\cdot (x_1-s)^2\]
is concave near the origin.
Since $\eps$ is small, we can assume that the $V_s$ are convex subsets of~$\EE^m$.

Further note that 
\[K'=\bigcup_{s\in[-\eps,\eps]}V_s.\]

Also, the same argument as in \ref{CAT(0)?}$C$ shows that
\begin{clm}{}\label{eq:VcVnV}
If $a<b<c$, then $V_b\supset V_a\cap V_c$.

\end{clm}

Given an array of values $s_1<\ldots<s_k$ in $[-\eps,\eps]$,
set $V_i=V_{s_i}$ and
consider the unions 
\[W_i=V_1\cup\ldots\cup V_i\]
equipped with the induced length metric.

Note that the array $(s_1,\ldots,s_k)$ can be chosen in such a way that
$W_k$ is arbitrarily close to~$K'$ in the sense of Hausdorff.

Arguing as in \ref{CAT(0)?}$C$, we get that the following claim implies \ref{K'-is-CAT}.
\begin{clm}{}
All $W_i$ are $\CAT(0)$.
\end{clm}

This claim is proved by induction.
Base: $W_1=V_1$ is $\CAT(0)$ as a convex subset in~$\EE^m$.

\parit{Induction step:} Assume that $W_i$ is $\CAT(0)$.
According to \ref{eq:VcVnV}, 
\[V_{i+1}\cap W_i=V_{i+1}\cap V_i.\]
Moreover, this is a convex set in $\EE^m$ 
and therefore it is a convex set in $W_i$ and in $V_{i+1}$.
By the Reshetnyak gluing theorem (\ref{thm:gluing}), $W_{i+1}$ is $\CAT(0)$.
Hence the claim follows.
\claimqeds

Summarizing, we have proved the following claim:

\begin{clm}{}\label{clm-strong2convex}
$K'$ is $\CAT(0)$ if $K$ is 
\index{strongly two-convex set}\emph{strongly two-convex},
that is, $k_2(p)>0$ at any point $p\in S$.
\end{clm}

It remains to show that $p$ admits a $\CAT(0)$ neighborhood in the case $k_2(p)\z=0$.

Choose a coordinate system $(x_1,\ldots,x_m)$ as above,
so that the $(x_1,\ldots,x_{m-1})$-coordinate hyperplane is the tangent subspace to $S$ at~$p$.

Fix $\eps>0$ so that a neighborhood of $p$ in $S$ 
is the graph
\[x_m= f(x_1,\ldots,x_{m-1})\]
of a function $f$ defined on the open ball $B$ of radius $\eps$ centered at the origin in the $(x_1,\ldots,x_{m-1})$-hyperplane.
Fix a smooth positive 
strongly 
convex function $\phi\:B\to \RR_+$
such that $\phi(x)\to\infty$ as $x$ approaches the boundary of $B$.
Note that for $\delta>0$, the subgraph $K_\delta$ defined by the inequality
\[x_m\le f(x_1,\ldots,x_{m-1})-\delta\cdot\phi(x_1,\ldots,x_{m-1})\]
is strongly two-convex.
By \ref{clm-strong2convex}, $K_\delta$ is $\CAT(0)$.

Finally as $\delta\to0$, the closed $\eps$-neighborhoods of $p$ in $K_\delta$ 
converge to the closed $\eps$-neighborhood of $p$ in $K$.
By \ref{prop:cat-limit}, the $\eps$-neighborhood of $p$ is $\CAT(0)$.
\qeds

Note that two-convexity survives under affine transformations of a Euclidean space.
Therefore, as a consequence of the smooth two-convexity theorem (\ref{thm:set-with-smooth-bry:CBA}), we get the following.

\begin{thm}{Corollary}\label{cor:affine}
Let $K\subset \EE^m$ be a closed connected set that admits a smooth parametrization by a closed $m$-ball.
If $K$ equipped with the induced length metric is $\CAT(0)$, then the same holds for any affine transformation of $K$.
\end{thm}

Corollary~\ref{cor:shefel} will provide a more general statement in the 3-dimensional Euclidean space.

\section{Open plane sets}

In this section, we consider inheritance of upper curvature bounds by subsets of the Euclidean plane.

\begin{thm}{Theorem}\label{thm:bishop-plane}
Let $\Omega$ be an open simply connected subset of~$\EE^2$.
Equip $\Omega$ with the induced length metric and denote its completion 
by~$K$.
Then $K$ is $\CAT(0)$.
\end{thm}

If $\Omega$ is not open, then the induced intrinsic metric might be infinite for some pair of points.
In this case the proof implies that each metric component of $\Omega$ with the induced length metric is $\CAT(0)$.

\parit{Sketch of proof.}
It is sufficient to show that any triangle in $K$ is thin,
as defined in \ref{def:k-thin}.

Note that $K$ admits a length-preserving map to $\EE^2$ that extends the embedding $\Omega\hookrightarrow\EE^2$.
Therefore each triangle $\trig xyz$ in $K$ can be mapped to the plane in a length-preserving way.
Since $\Omega$ is simply connected, any open region, say $\Delta$, that is surrounded by the image of $\trig xyz$ lies completely in~$\Omega$.

\begin{figure}[b!]
\vskip-0mm
\centering
\includegraphics{mppics/pic-1203}
\caption{A triangle with overlapping sides.}
\label{pic-1203}
\end{figure}

Note that in each triangle $\trig xyz$ in $K$, the sides $[xy]$, $[yz]$ and $[zx]$ intersect each other along a geodesic starting at a common vertex, possibly a one-point geodesic.
In other words, every triangle in $K$ looks like the one in figure \ref{pic-1203}.

Indeed, assuming the contrary, there will be a {}\emph{lune} in $K$ bounded by two minimizing geodesics with common ends but no other common points.
The image of this lune in the plane must have concave sides, since otherwise one could shorten the sides by pushing them into the interior.
Evidently, there is no plane lune with concave sides ---  a contradiction.

Note that it is sufficient to consider only simple triangles $\trig xyz$, 
that is, triangles whose sides $[xy]$, $[yz]$ and $[zx]$ intersect each other only at the common vertices.
If this is not the case, chopping the overlapping part of sides reduces to the case when the triangle is simple (this is formally stated in Exercise~\ref{ex:chopping-triangle}).

Again, the open region, say $\Delta$, bounded by the image of $\trig xyz$ has concave sides in the plane, since otherwise one could shorten the sides by pushing them into~$\Omega$.
It remains to solve Exercise~\ref{ex:concave-triangle}.
\qeds

\begin{thm}[!]{Exercise}\label{ex:chopping-triangle}
Assume that $[pq]$ is a common part of the two sides $[px]$ and $[py]$ of the triangle $[pxy]$.
Consider the triangle $[qxy]$ whose sides are formed by arcs of the sides of~$[pxy]$.
Show that if $[qxy]$ is thin, then so is~$[pxy]$.
\end{thm}

\begin{thm}[!]{Exercise}\label{ex:concave-triangle}
Assume $S$ is a closed plane 
region whose boundary is a plane triangle $T$ with concave sides.
Equip $S$ with the induced length metric.
Show that the triangle $T$ is thin in~$S$.
\end{thm}

\begin{thm}{Proposition}\label{prop:bishop-sphere}
Let $\Theta$ be an open connected subset of the unit sphere $\mathbb{S}^2$ that does not contain a closed hemisphere.
Equip $\Theta$ with the induced length metric.
Let $\tilde \Theta$ be a metric cover of $\Theta$ 
such that any closed curve in $\tilde \Theta$ shorter than $2\cdot\pi$ is contractible.

Show that the completion of $\tilde \Theta$ is $\CAT(1)$.
\end{thm}


This proposition will be used in the next section.
It is a spherical analog of Theorem \ref{thm:bishop-plane}
and can be proved along the same lines with an additional idea from the following exercise.

\begin{thm}{Exercise}\label{ex:bishop-sphere}
Let $K$ be a closed subset of the unit sphere $\mathbb{S}^2$.
Suppose $K$ is simply connected, bounded by a simple closed Lipschitz curve, and its interior does not contain a closed hemisphere.
Show that $K$ with the induced length metric is $\CAT(1)$.
\end{thm}

\section{Two-convexity theorem}

\begin{thm}{Theorem}\label{thm:shefel}
Let $\Omega$ be a connected open set in~$\EE^3$.
Equip $\Omega$ with the induced length metric
and denote by $\tilde K$ the completion of the universal metric cover of~$\Omega$.
Then $\tilde K$ is $\CAT(0)$ 
if and only if $\Omega$ is two-convex.
\end{thm}

The proof of this statement will be given in the following three sections.
First we prove its polyhedral analog, then we prove some properties of two-convex hulls in the 3-dimensional Euclidean space and only then do we prove the general statement.

The following exercise shows that the analogous statement does not hold in higher dimensions.

\begin{thm}{Exercise}\label{ex:two-planes}
Let $\Pi_1,\Pi_2$ be two planes in $\EE^4$ intersecting at a single point.
Let $\tilde K$ be the completion of the universal metric cover of $\EE^4\setminus(\Pi_1\cup\Pi_2)$.

Show that 
$\tilde K$ is $\CAT(0)$ if and only if $\Pi_1\perp\Pi_2$.

\end{thm}

Before coming to the proof of the two-convexity theorem, 
let us formulate a few corollaries.
The following corollary is a generalization of the smooth two-convexity theorem (\ref{thm:set-with-smooth-bry:CBA}) for 3-dimensional Euclidean space.

\begin{thm}{Corollary}\label{cor:shefel}
Let $K$ be a closed subset in $\EE^3$.
Assume there is a Lipschitz homeomorphism from a closed ball in $\EE^3$ to $K$.
Then $K$ with the induced length metric is $\CAT(0)$ if and only if the interior of $K$ is two-convex.
\end{thm}

\parit{Proof.}
Set $\Omega=\Int K$.
By the domain invariance theorem, $\Omega$ is homeomorphic to an open ball; in particular it is simply connected.

Apply the two-convexity theorem to~$\Omega$.
Note that the completion of $\Omega$ equipped with the induced length metric 
is isometric to $K$ with the induced length metric.
Hence the result.
\qeds

Note that the Lipschitz condition is used just once to show that the completion of $\Omega$ is isometric to $K$ with the induced length metric.
This property holds for a wider class of hypersurfaces;
for instance, some examples of Alexander horned balls have $\CAT(0)$ induced length metric.

Let $U$ be an open set in~$\RR^2$.
A continuous function $f\:U\to\RR$ is called 
\index{saddle function}\emph{saddle} 
if for any linear function $\ell\:\RR^2\to\RR$, the difference 
$f-\ell$
does not have local maxima or local minima in~$U$.
Equivalently, the open subgraph and epigraph of $f$
\begin{align*}
&\set{(x,y,z)\in\EE^3}{z<f(x,y),\ (x,y)\in U},
\\
&\set{(x,y,z)\in\EE^3}{z>f(x,y),\ (x,y)\in U}
\end{align*}
are two-convex. 

\begin{thm}{Theorem}\label{thm:shefel-graph}
Let $f\:\DD\to \RR$ be a Lipschitz function which is saddle in the interior of the closed unit disc~$\DD$. 
Then the graph
\[\Gamma=\set{(x,y,z)\in \EE^3}{z=f(x,y)},\] 
equipped with the induced length metric is $\CAT(0)$.
\end{thm}

\parit{Proof.}
Since the function $f$ is Lipschitz,
its graph $\Gamma$ with the induced length metric is bi-Lipschitz equivalent to $\DD$ with the Euclidean metric.

\begin{figure}[t!]
\vskip-0mm
\centering
\includegraphics{mppics/pic-1206}
\caption{}
\end{figure}

Consider the sequence of sets 
\[K_n
=
\set{(x,y,z)\in \EE^3}{z\lege f(x,y)\pm\tfrac1n,\ (x,y)\in \DD}.\]
Note that each $K_n$ meets the conditions in Corollary~\ref{cor:shefel}.
Therefore, $K_n$ equipped with the induced length metric is $\CAT(0)$.
It remains to note that $K_n\to \Gamma$ in the sense of Gromov--Hausdorff, and apply \ref{prop:cat-limit}.
\qeds

\section{Polyhedral case}

Now we are back to the proof of the two-convexity theorem (\ref{thm:shefel}).

Recall that a subset $P$ of $\EE^m$ is called \index{polyhedral set}\emph{polyhedral}
if it can be presented as a union of a finite number of simplices.
Similarly,
a \index{spherical polyhedral set}\emph{spherical polyhedral set}
is a union of a finite number of simplices in~$\mathbb{S}^m$.

Note that any polyhedral set admits a finite triangulation.
Therefore any polyhedral set equipped with the induced length metric
forms a Euclidean polyhedral space as defined in~\ref{def:poly}.
 
\begin{thm}{Lemma}\label{lem:poly-shefel}
The two-convexity theorem (\ref{thm:shefel}) holds if the set $\Omega$ is the interior of a polyhedral set.
\end{thm}

The statement might look obvious, but we will face a technical difficulty in the proof that is related to the following.
Let $P$ be a polyhedral set and $\Omega$ its interior,
both considered with the induced length metrics.
If the completion $K$ of $\Omega$
is isometric to~$P$, then the lemma follows easily from
\ref{thm:PL-CAT}.

\begin{figure}[b!]
\vskip-0mm
\centering
\includegraphics{mppics/pic-1207}
\caption{A polyhedral set that is not isometric to the completion of its interior.}
\label{pic-1207}
\end{figure}

However in general
we only have a locally distance-preserving map $K\to P$;
it does not have to be onto and it may not be injective. 
An example can be guessed from the figure~\ref{pic-1207}.
Nevertheless, it is easy to see that $K$ is always a polyhedral space.

The proof uses the following two exercises.

\begin{thm}[!]{Exercise}\label{ex:hemisphere}
Show that any closed path of length $<2\cdot \pi$ in the unit sphere $\mathbb{S}^2$ lies in an open hemisphere.
\end{thm}

\begin{thm}[!]{Exercise}\label{ex:inner-support}
Assume $\Omega$ is an open subset in $\EE^3$
that is not two-convex.
Show that there is a plane $W$ such that the complement 
$W\setminus \Omega$ contains an isolated point and a small circle around this point in $W$ is contractible in $\Omega$.
\end{thm}

\parit{Proof of \ref{lem:poly-shefel}.}
The ``only if'' part can be proved in the same way as in the smooth two-convexity theorem (\ref{thm:set-with-smooth-bry:CBA}) with additional use of Exercise~\ref{ex:inner-support}.

\parit{If part.}
Assume that $\Omega$ is two-convex.
Denote by $\tilde\Omega$ the universal metric cover of~$\Omega$.
Let $\tilde K$ and $K$ be the completions of $\tilde\Omega$ and $\Omega$, respectively.

The main step is to show that $\tilde K$ is $\CAT(0)$. 

Note that $K$ is a polyhedral space and the covering $\tilde\Omega\to\Omega$ extends to a covering map $\tilde K\to K$ which might be branching at some vertices.%
\footnote{For example, if $K=\set{(x,y,z)\in\EE^3}{|z|\le |x|+|y|\le 1}$ and $p$ is the origin, then $\Sigma_p$,
the space of directions at $p$,
is not simply connected and $\tilde K\to K$ branches at~$p$.}

Fix a point $\tilde p\in \tilde K\setminus\tilde\Omega$; 
denote by $p$ the image of $\tilde p$ in~$K$.
Note that $\tilde K$ is a ramified cover of $K$ and hence is locally contractible.
Thus, any loop in $\tilde K$ is homotopic to a loop in $\tilde\Omega$.
Since $\tilde\Omega$ is simply connected, so is $\tilde K$.

Thus, by the globalization theorem (\ref{thm:hadamard-cartan}), it is sufficient to show that

\begin{clm}{}\label{eq:curv=<0}
a small neighborhood of $\tilde p$ in $\tilde K$ is $\CAT(0)$.
\end{clm}

Recall that $\Sigma_{\tilde p}=\Sigma_{\tilde p}\tilde K$ denotes the space of directions at~$\tilde p$.
Note that a small neighborhood of $\tilde p$ in $\tilde K$
is isometric to an open set in the cone over $\Sigma_{\tilde p}\tilde K$.
By \ref{ex:cone+susp}, \ref{eq:curv=<0} follows once we can show that

\begin{clm}{}\label{eq:curv=<1}
$\Sigma_{\tilde p}$ is $\CAT(1)$.
\end{clm}

By rescaling, we can assume that all faces of $K$ that do not contain $p$ lie at distance at least 2 from~$p$.
Let $\mathbb{S}^2$ be the unit sphere centered at $p$,
and let $\Theta=\mathbb{S}^2\cap\Omega$.
Note that $\Sigma_pK$ is isometric to the completion of $\Theta$
and $\Sigma_{\tilde p}\tilde K$ is the completion of the regular metric covering $\tilde\Theta$ of $\Theta$ induced by the universal metric cover $\tilde \Omega\to \Omega$.

By \ref{prop:bishop-sphere}, it remains to show the following:
\begin{clm}{}
Any closed curve in $\tilde\Theta$ shorter than $2\cdot\pi$ is contractible.
\end{clm}

Fix a closed curve $\tilde \gamma$ of length $<2\cdot\pi$ in~$\tilde\Theta$.
Its projection $\gamma$ in $\Theta\subset\mathbb{S}^2$ has the same length.
Therefore, by Exercise~\ref{ex:hemisphere}, $\gamma$ lies in an open hemisphere.
Then for a plane $\Pi$ passing close to $p$,
the central projection $\gamma'$ of $\gamma$ to $\Pi$ is defined and lies in~$\Omega$.
By construction of $\tilde\Theta$, the curve $\gamma$ and therefore $\gamma'$ are contractible in~$\Omega$.
From two-convexity of $\Omega$
and Proposition~\ref{prop:3d-strong-2-convexity},
the curve $\gamma'$ is contractible in $\Pi\cap \Omega$.

It follows that $\gamma$ is contractible in $\Theta$ 
and therefore $\tilde\gamma$ is contractible in~$\tilde\Theta$.
\qeds

\section{Two-convex hulls}\label{sec:Two-convex hulls}

The following proposition 
describes a construction which produces the two-convex hull $\Conv_2 \Omega$ of an open set $\Omega\subset\EE^3$.
This construction is very close to the one given by Samuel Shefel~\cite{shefel-1964}.

\begin{thm}{Proposition}\label{prop:2-conv-construction}
Let $\Pi_1,\Pi_2\ldots$ be an everywhere dense
sequence of planes in~$\EE^3$.
Given an open set $\Omega$, consider 
the recursively defined sequence of open sets 
$\Omega=\Omega_0\subset\Omega_1\subset\ldots$
such that 
$\Omega_n$ is the union of $\Omega_{n-1}$ 
and all the bounded components of 
$\EE^3\setminus(\Pi_n\cup \Omega_{n-1})$.
Then 
\[\Conv_2\Omega=\bigcup_n\Omega_n.\]

\end{thm}

\parit{Proof.}
Set 
\[\Omega'=\bigcup_n\Omega_n.\eqlbl{eq:Omega'}\]
Note that $\Omega'$ is a union of open sets; in particular, $\Omega'$ is open.

Let us show that
\[\Conv_2\Omega\supset\Omega'.\eqlbl{eq:Conv2>Omega'}\]
Suppose we already know that $\Conv_2\Omega\supset\Omega_{n-1}$. 
Fix a bounded component $\mathfrak{C}$ of $\EE^3\setminus(\Pi_n\cup \Omega_{n-1})$.
It is sufficient to show that $\mathfrak{C}\subset\Conv_2\Omega$.

By \ref{prop:two-hull-open}, $\Conv_2\Omega$ is open.
Therefore, if $\mathfrak{C}\not\subset\Conv_2\Omega$,
then there is a point $p\in \mathfrak{C}\setminus\Conv_2\Omega$ lying at maximal distance from~$\Pi_n$.
Denote by $W_p$ the plane containing $p$ which is parallel to~$\Pi_n$.

Note that $p$ lies in a bounded component of $W_p\setminus \Conv_2\Omega$.
In particular $p$ can be surrounded by a simple closed curve $\gamma$ in $W_p\z\cap\Conv_2\Omega$.
Since $p$ lies at maximal distance from $\Pi_n$,
the curve $\gamma$ is null-homotopic in $\Conv_2\Omega$.
Therefore $p\in \Conv_2\Omega$ ---  a contradiction.

By induction, $\Conv_2\Omega\supset\Omega_n$ for each~$n$.
Therefore \ref{eq:Omega'} implies~\ref{eq:Conv2>Omega'}.

It remains to show that $\Omega'$ is two-convex.
Assume the contrary; 
that is, there is a plane $\Pi$ 
and a simple closed curve $\gamma\:\mathbb{S}^1\to \Pi\cap \Omega'$ 
which is null-homotopic in $\Omega'$,
but not null-homotopic in $\Pi\cap\Omega'$.

By approximation we can assume that $\Pi=\Pi_n$ for a large $n$, and that $\gamma$ lies in $\Omega_{n-1}$.
By the same argument as in the proof of Proposition~\ref{prop:3d-strong-2-convexity} using the loop theorem, we can assume that there is an embedding $\phi\: \DD\to \Omega'$ such that $\phi|_{\partial\DD}=\gamma$ and $\phi(D)$ lies entirely in one of the half-spaces bounded by~$\Pi$.
By the $n$-th step of the construction, the entire bounded domain $U$ bounded by $ \Pi_n$ and $\phi(D)$ is contained in $\Omega'$ and hence $\gamma$ is contractible in $\Pi\cap\Omega'$ --- a contradiction.
\qeds

\begin{thm}{Key lemma}\label{lem:key-shefel}
The two-convex hull of the interior of a polyhedral set in $\EE^3$
is also the interior of a polyhedral set.
\end{thm}

\parit{Proof.}
Fix a polyhedral set $P$ in~$\EE^3$.
Set $\Omega=\Int P$.
We may assume that $\Omega$ is dense in~$P$
(if not, redefine $P$ as the closure of $\Omega$).
Denote by $F_1,\ldots,F_m$ the facets of~$P$.
By subdividing $F_i$ if necessary, we may assume that all $F_i$ are convex polygons.

Set $\Omega'=\Conv_2\Omega$ and let $P'$ be the closure of~$\Omega'$.
Further, let $F'_i=F_i\setminus \Omega'$;
in other words,
$F'_i$ is the subset of the facet $F_i$ 
which remains on the boundary of~$P'$.

From the construction of the two-convex hull (\ref{prop:2-conv-construction}) we have:

\begin{clm}{}\label{clm:F'-convex}
$F'_i$ is a convex subset of~$F_i$.
\end{clm}

Further, since $\Omega'$ is two-convex we obtain the following:

\begin{clm}{}\label{clm:complement-of-F'-convex}
Each connected component of the complement $F_i\setminus F'_i$ is convex.
\end{clm}

\begin{figure}[t!]
\vskip-0mm
\centering
\includegraphics{mppics/pic-1209}
\caption{Facet $F_i$ and its subset $F'_i$ which remains on the boundary of~$P'$.}
\end{figure}

Indeed, assume a connected component $A$ of $F_i\setminus F'_i$ fails to be convex.
Then $F'_i$ has an extreme point in the interior of $F_i$.
In other words, there is a supporting line $\ell$ to $F'_i$ touching $F'_i$ at a single point in the interior of~$F_i$.
Then one could rotate the plane of $F_i$ slightly around $\ell$ and move it parallelly to cut a ``cap'' from the complement of~$\Omega$.
The latter means that $\Omega$ is not two-convex
--- a contradiction.
\claimqeds


From \ref{clm:F'-convex} and \ref{clm:complement-of-F'-convex}, we conclude
\begin{clm}{}\label{clm:F'} $F'_i$ is a convex polygon for each~$i$.
\end{clm}

Consider the complement 
$\EE^3\setminus \Omega$ 
equipped with the length metric.
By construction of the two-convex hull (\ref{prop:2-conv-construction}), 
the complement $L\z=\EE^3\setminus (\Omega'\cup P)$
is locally convex;
that is, any point of $L$ admits a convex neighborhood.

Summarizing: (1)
$\Omega'$ is a two-convex open set,
(2) the boundary $\partial\Omega'$ 
contains a finite number of polygons $F_i'$
and the remaining part $S$ of the boundary is locally concave.
It remains to do the following exercise.

\begin{thm}[!]{Advanced exercise}\label{ex:convex+saddle+broken=>PL}
Let $\bar S$ be the closure of $S$.
Denote by $\partial S$ the boundary of $\bar S$ in $\partial \Omega'$.

\begin{subthm}{ex:convex+saddle+broken=>PL:a}
Show that (1) and (2) imply that any point $p\in S$ lies in a line segment in $\bar S$ with ends in $\partial S$ or in a polygon in $\bar S$ with vertices in~$\partial S$.
\end{subthm}

\begin{subthm}{ex:convex+saddle+broken=>PL:b}
Use part \ref{SHORT.ex:convex+saddle+broken=>PL:a} to show that $\bar S$ (and therefore $\partial\Omega'$)
is piecewise linear. \qeds
\end{subthm}

\end{thm}

\section{Proof of Shefel's theorem}

\parit{Proof of \ref{thm:shefel}.}
The ``only if'' part can be proved in the same way as in the smooth two-convexity theorem (\ref{thm:set-with-smooth-bry:CBA}) with the additional use of Exercise~\ref{ex:inner-support}.

\parit{If part.}
Suppose $\Omega$ is two-convex. 
We need to show that $\tilde K$ is $\CAT(0)$.

Fix a quadruple of points $x_1,x_2,x_3,x_4\in\tilde \Omega$.
Let us show that
$\CAT(0)$ comparison holds for this quadruple.

Fix $\eps>0$.
Choose six polygonal lines in $\tilde \Omega$ connecting all pairs of points $x_1,x_2,x_3,x_4$, where the length of each polygonal line is at most $\eps$ bigger than
the distance between its ends in the length metric on~$\tilde \Omega$.
Denote by $X$ the union of these polygonal lines.
Choose a polyhedral set $P$ in $\Omega$ such that its interior $\Int P$ contains the projections of all six polygonal lines and discs which contract all the loops created by them (it is sufficient to take 3 discs).

Denote by $\Omega'$ the two-convex hull of the interior of~$P$.
According to the key lemma (\ref{lem:key-shefel}), $\Omega'$ is the interior of a polyhedral set.

Equip $\Omega'$ with the induced length metric.
Consider the universal metric cover $\tilde\Omega'$ of~$\Omega'$.
(The covering $\tilde\Omega'\to\Omega'$ might be nontrivial ---
even if $\Int P$ is simply connected, its two-convex hull $\Omega'$ might not be simply connected.)
Denote by $\tilde K'$ the completion of~$\tilde\Omega'$.

By Lemma~\ref{lem:poly-shefel}, $\tilde K'$ is $\CAT(0)$.

By construction of $\Int P$, the embedding $\Int P\hookrightarrow \Omega'$
admits a lift $\iota\:X\hookrightarrow \tilde K'$, and $\iota$ almost preserves the distances between the points $x_1,x_2,x_3,x_4$;
namely 
\begin{align*}
\dist{\iota(x_i)}{\iota(x_j)}{\tilde K'}\lessgtr \dist{x_i}{x_j}{\Int P}\pm\eps.
\end{align*}

Since $\eps>0$ is arbitrary and $\CAT(0)$ comparison holds in $\tilde K'$,
we get that $\CAT(0)$ comparison holds for $x_1,x_2,x_3,x_4$ in $\tilde \Omega$.

The statement follows since the quadruple $x_1,x_2,x_3,x_4\in\tilde\Omega$ is arbitrary.
\qeds

\begin{thm}{Exercise}\label{ex:CAT=>two-convex}\label{ex:two-convex-not-a-CAT}
Assume $K\subset \EE^{m}$ is a subgraph of a Lipschitz function $f\:\RR^{m-1}\to\RR$.
Equip $K$ with the induced length metric.
Show that if $K$ is $\CAT(0)$, then $K$ is two-convex.

Show that the converse does not hold.
\end{thm}

\section{Notes}

Under the name {}\emph{$(n-2)$-convex sets}, 
two-convex sets in $\EE^n$ were introduced by Mikhael Gromov \cite{gromov-1991}.
In addition to the inheritance of upper curvature bounds by two-convex sets discussed in this lecture, 
these sets appear as the maximal open sets with vanishing curvature in Riemannian manifolds with non-negative or non-positive sectional curvature [see Lemma 5.8 in \ncite{buyalo}, \ncite{andersson-howard}
and \ncite{panov-petrunin}].

Two-convex sets could be defined using homology instead of homotopy, as in Gromov's formulation of the Lefschetz theorem \cite[\S\textonehalf]{gromov-1991}.
Namely, we can say that $K$ is two-convex if the following condition holds: \textit{if a 1-dimensional cycle $z$ has support in the intersection of $K$ with a plane $W$ and bounds in $K$, then it bounds in $K\cap W$}.

The resulting definition is equivalent to the one used above.
But unlike our definition it can be generalized to define $k$-convex sets in $\EE^m$ for~$k>2$.
With this homological definition one can also avoid the use of the loop theorem, whose proof is quite involved.
We stick to the homotopy-based definition since it is slightly easier to visualize.

Both definitions work well for open sets; for general sets one should be able to give a similar definition using  an appropriate homotopy/ho\-mo\-logy theory.

Proposition~\ref{prop:two-cove+smooth} was proved by Mikhael Gromov \cite[\S\textonehalf]{gromov-1991}, but we added a few details.

Jointly with David Berg and Richard Bishop, the first author proved the following theorem, which gives the exact upper bound on Alexandrov's curvature for the Riemannian manifolds with boundary.
This theorem includes the smooth two-convexity theorem (\ref{thm:set-with-smooth-bry:CBA}) as a partial case.

\begin{thm}{Theorem}
Let $M$ be a Riemannian manifold with boundary~$\partial M$.
A direction tangent to the boundary will be called concave if there is a short geodesic in this direction which leaves the boundary and goes into the interior of~$M$.
A sectional direction (that is, a 2-plane) 
tangent to the boundary 
will be called {}\emph{concave} if all the directions in it are concave.

Denote by $\kappa$ an upper bound of sectional curvatures of $M$ and  
sectional curvatures of $\partial M$ in the concave sectional directions. 
Then $M$ is locally $\CAT(\kappa)$. 
\end{thm}

\begin{thm}{Corollary}
Let $M$ be a Riemannian manifold with boundary~$\partial M$. 
Assume that all the sectional curvatures of $M$ and $\partial M$ are bounded above by~$\kappa$.
Then $M$ is locally $\CAT(\kappa)$.
\end{thm}

The Theorem~\ref{thm:bishop-plane} was proved by Richard Bishop \cite{bishop}.
Generalizations of this result were proved by Alexander Lytchak, Stephan Stadler and Stefan Wenger; see \cite{lytchak-stadler} and \cite[Proposition 12.1]{lytchak-wenger}.

Samuel Shefel was intrigued by the survival of metric properties under affine transformation.
An instance of this can be seen in Corollary~\ref{cor:affine}.

Theorem~\ref{thm:shefel-graph} is from Shefel's original paper~\cite{shefel-1965}.
It is related to Alexandrov's theorem about ruled surfaces \cite{alexandrov-1957-ruled-surfaces}.

The following question remains open.

\begin{thm}{Shefel's conjecture}
Suppose $s\:\DD \to\EE^3$ is a Lipschitz embedding of the disc $\DD$,
and its image $D=s(\DD)$ is \index{saddle surface}\emph{saddle}; that is, for any plane $\Pi$ each connected component of $D\setminus \Pi$ meets the boundary of $D$.
Then $D$ equipped with the length-metric is locally $\CAT(0)$.
\end{thm}

A partial answer is given in Theorem \ref{thm:shefel-graph},
which is a result of Shefel~\cite{shefel-1964}.
He also solved an analogous question in $\EE^2$; see \cite{shefel-1965}.
Applying Crofton-type argument to his result, one can see that a saddle disc satisfies the isoperimetric inequality $a\le C\cdot \ell^2$
where $a$ is the area of a disc bounded by a curve of length $\ell$ and $C=\tfrac{1}{3\cdot\pi}$.
By a result of Alexander Lytchak and Stefan Wenger \cite{lytchak-wenger}, Shefel's conjecture would follow if the constant could be made optimal; that is, $C=\tfrac{1}{4\cdot\pi}$.
Several variations and generalizations of this conjecture are discussed by the third author and Stephan Stadler \cite{petrunin-stadler}.
